% status: ZERO
% Chapter 12 of the second edition. Plan: docs/STRUCTURE.md. Rules: AUTHORING.md.
% Measurements: research/measurements/2026-09-23-m1-part2.md
\chapter{Quantizing the KV Cache}\label{ch:kv-quantization}

\begin{ChapterIntent}
At long context and large batch the KV cache, rather than the weights,
dominates both memory and the bytes read per step, and quantizing it is the
most direct relief. Its errors behave differently from weight errors: they
enter every attention score of every later token, and keys carry outlier
channels that per-token scaling handles badly. This chapter derives
per-channel key and per-token value quantization, FP8 caches with calibrated
scales, four-bit and lower schemes, and pre-RoPE and rotated variants, and it
shows how to measure long-context damage rather than trusting short-context
perplexity.
\end{ChapterIntent}

\OpeningSection{Twice the context in the same memory}

llama.cpp stores its KV cache in 16-bit floats unless told otherwise; one
flag stores it in the 8- or 4-bit block formats it uses for weights. For
Llama~3.2~3B the choice shrinks a token's keys and values from 112~KiB to
59.5~KiB at 8 bits or 31.5~KiB at 4 (derived from the formats: 34 and 18 bytes
per block of 32 values). A 16,384-token context whose cache needs 1.88~GB in
16 bits needs 1.00~GB at 8 and 0.53~GB at 4: in the same memory, 3.6 times as
many such conversations.

The speed does not follow the bytes. On the M1, with flash attention on and
16,384 tokens in the cache, llama.cpp decoded at \textbf{12.4} tokens per second
with the 16-bit cache, \textbf{9.5} with the 8-bit cache and \textbf{11.4} with
the 4-bit cache (measured; record \texttt{2026-09-23-m1-part2}). The quantized
caches were slower at every depth measured, from an empty cache to 16,384
tokens. Subtract each format's empty-cache step time and the attention over a
16,384-token cache read its bytes at about 46~GB/s in 16 bits, 16~GB/s at 8
bits and 14~GB/s at 4 (derived). In this engine, on this GPU, unpacking the
cache inside the kernel costs more time than the bytes it saves.

So the case for a quantized cache on this machine is capacity alone. The case
for speed needs two more things: a workload in which the cache is most of the
bytes, and a kernel that reads the smaller format at close to full bandwidth.
The napkin math below prices the first; only a measurement settles the second.

\NapkinSection{Napkin math: when the cache is most of the bytes}

For $B$ sequences with mean context $\bar{S}$, the cache's share of a decode
step's bytes is
\begin{equation}
\phi_{kv} = \frac{B\,\bar{S}\,m_{kv}}{W + B\,\bar{S}\,m_{kv}},
\label{eq:kv-share}
\end{equation}
and \Chref{napkin-math} found $\phi_{kv}\approx0.8$ for Qwen3-8B on an H100 with
a full memory of 4,096-token sequences. When the cache is most of the bytes,
shrinking it pays twice: every step reads less, and twice as many sequences
fit. Quantizing only the cache of that configuration from BF16 to FP8 --- the
weights untouched --- lets the H100 hold 197 sequences instead of 98 and, with
the step model of equation~\eqref{eq:step-time}, roughly doubles throughput
and halves the cost per token (derived, with the same illustrative
efficiencies, and assuming the kernel reads FP8 as efficiently as BF16 --- the
assumption the opening measurement warns about). No weight format change in
\Chref{weight-quantization} buys that much at this operating point.

\section{Why cache errors are different}

A quantized weight's error enters one layer's output for the current token. A
quantized key's error enters the attention score of \emph{every later query}
that looks at that token, for the rest of the request. And it enters through an
exponential: a score error $\Delta$ multiplies the token's attention weight by
$e^{\Delta}$ before normalization. An error that is harmless at 500 tokens can
accumulate into a wrong answer at 50,000, because the model's ability to
retrieve one fact from far back depends on one key standing out from many.
Values are gentler: their errors are averaged by the attention weights.

\section{Keys have outlier channels}

Key vectors are not evenly distributed across their dimensions: a few channels
carry consistently large magnitudes, the same channels token after token. KIVI
measured this and drew the design conclusion: quantize keys \emph{per channel},
so that each channel gets its own scale across a group of tokens, and values
\emph{per token}, because their outliers are not channel-aligned
\cite{liu2024kivi}. With that split, KIVI stored both at 2 bits with little
quality loss and fit up to four times larger batches.

\section{Per-channel keys, per-token values}

The asymmetry has a systems cost. A per-token scale can be computed the moment a
token's key is produced. A per-channel scale spans many tokens, so it cannot be
fixed until a group of tokens exists; recent tokens wait in higher precision
until their group is complete and can be quantized together. The engine's
cache becomes two-tiered --- a quantized body and a small full-precision tail
--- and every attention kernel must read both.

\section{FP8 caches and their scales}

The simplest production option stores the cache in FP8. vLLM offers E4M3 and
E5M2 caches, with one scale per tensor or per attention head, computed from a
calibration dataset or left at a default of 1.0 \cite{vllm2026kvcache}.
FP8 halves the cache with small, usually negligible error --- provided the
scale is right. E4M3's largest value is 448 \cite{micikevicius2022fp8}; a cache
whose values exceed the range at a scale of 1.0 saturates, and one whose values
are tiny wastes most of the format's precision. ``Usually negligible'' is a
claim to measure per model.

\section{Four bits and below}

Below eight bits the cache needs the channel-aware treatment above, or
rotations. KVQuant combined per-channel keys, pre-RoPE quantization,
non-uniform per-layer code books and separately stored outliers, and reported
under 0.1 perplexity increase at 3 bits --- enough to serve LLaMA-7B with a
million-token context on one 80~GB GPU \cite{hooper2024kvquant}. Rotation
methods spread outliers across channels before quantizing: QuaRot quantizes the
cache to 4 bits alongside weights and activations \cite{ashkboos2024quarot}.
TurboQuant rotates each vector randomly and applies an optimal scalar quantizer
per coordinate, and reports quality neutrality at 3.5 bits per channel
\cite{zandieh2025turboquant}; vLLM added it as an attention backend in April
2026, with FP8 keys and 4-bit values as its most conservative preset
\cite{vllm2026turboquantpr}.

\section{Before or after RoPE; rotations}

Rotary position embedding rotates pairs of key channels by angles that depend
on position (Appendix~\ref{app:rope}), which mixes the outlier channels with
their partners and makes their statistics depend on position. Quantizing keys
\emph{before} RoPE preserves the channel structure that per-channel scales
exploit, at the price of applying RoPE when the key is read rather than when it
is written \cite{hooper2024kvquant}. Rotation methods use the opposite idea:
apply a known orthogonal rotation that spreads outliers evenly, so that plain
per-token quantization works.

\section{Kernels that read a quantized cache}

A quantized cache is only useful if the attention kernel reads it without
first expanding it back to 16 bits in memory, which would spend the bandwidth it
saved. The kernel dequantizes blocks as it loads them, inside the tiled loop of
\Chref{flashattention} and \Chref{decode-attention}. llama.cpp makes the
coupling explicit: it refuses a quantized value cache unless flash attention is
enabled (llama.cpp \texttt{src/llama-context.cpp} at commit \texttt{389ff61d}),
tying the cache's format to the fused kernel. The dequantization is extra
arithmetic in what should be a bandwidth-bound loop, and it is not always
free. In the opening measurement, llama.cpp's Metal kernels read an 8-bit cache
at about a third of the rate at which they read a 16-bit one, so the smaller
cache was the slower one. Whether a format pays in time is a property of the
kernel and the hardware, not of the format.

\section{Measuring long-context damage}

Perplexity on short passages measures almost nothing about a cache: the
errors that matter accumulate over long contexts and show up as failures to
retrieve or use distant information. Measure with long inputs: retrieval of
planted facts across the whole context, long-document question answering, and
comparisons of the quantized and unquantized models' output distributions at
many context lengths. The TurboQuant integration in vLLM reported its FP8-key,
4-bit-value preset keeping 100\% on a needle-in-a-haystack test while
accuracy on a mathematics benchmark fell from 0.900 to about 0.860 on
Qwen3-4B \cite{vllm2026turboquantpr} --- one metric unmoved, another clearly
moved, by the same change.

\LedgerSection

Storing the KV cache in 8 or 4 bits instead of 16:

\begin{ConstraintLedger}
Memory capacity & buys & 1.88$\times$ the tokens per byte at 8 bits and 3.56$\times$ at 4 (llama.cpp's q8\_0 and q4\_0 caches: 8.5 and 4.5 bits per value). \\
Memory bandwidth & buys & Fewer cache bytes per step (derived) --- time saved only if the kernel reads the smaller format at full speed; on the M1, llama.cpp's did not, and the 16-bit cache was fastest at every measured depth. \\
Compute & spends & Dequantization inside the attention kernel: on the M1, enough to cut the cache's effective read rate from about 46 to 14--16~GB/s (derived). \\
Correctness & risks & Errors in every later attention score; measure at long context, not with short perplexity. \\
\end{ConstraintLedger}

\LabSection{quantize a cache and watch attention drift}

\LabIntent{In \texttt{code/labs/ch12-kv-quant/}, a Rust program that records
real keys and values from a small model's forward pass (or synthetic ones with
planted outlier channels), quantizes them per token and per channel at 8, 4 and
2 bits, and measures the error in attention outputs as the context grows,
against the full-precision oracle. The reader finds the context length at which
per-token 4-bit keys first change which token receives the most attention.}

\HappenedSection{half the bytes, slower}

When Hermon moved its KV-cache writes into C, storing the cache in 16-bit
floats was \textbf{slower} than storing it in 32-bit: 6.39~ms against 1.16~ms
for the same bulk write, despite moving half the bytes. The conversion from
32- to 16-bit floats was written in software with five branches per element
--- for NaN, overflow and subnormals --- and it cost more than the bytes it
saved while preventing the loop from vectorizing. Replacing it with the ARM
instruction that converts four values at once made the 16-bit write 8.8 times
faster and, as it should always have been, faster than the 32-bit one (Hermon
\texttt{docs/KERNEL\_DESIGN.md}, section 1; the kernels are selected by an
environment variable, so PREVIEW). A smaller format is a promise about bytes,
not about time; the conversion is part of the cost.

\FrontierSection{2026-09-24}

KV-cache formats moved from research into engines' command lines between 2024
and 2026: FP8 caches are routine \cite{vllm2026attentionbackends}, and
rotation-plus-scalar schemes such as TurboQuant ship behind a flag
\cite{vllm2026turboquantpr}. Architectures attack the same bytes from the other
side --- latent attention stores a compressed cache by design
(\Chref{shrinking-kv}), and DeepSeek-V4 compresses along the sequence, reporting
a tenth of DeepSeek-V3.2's cache at a million tokens \cite{deepseek2026v4}.
\todo{verify: NVFP4 KV-cache support in NVIDIA's serving stack from a primary
source}.

\ExercisesSection

\begin{enumerate}
\item \emph{Derivation.} Using equation~\eqref{eq:kv-share}, find the context
length at which quantizing the cache to FP8 speeds a batch of 32 Qwen3-8B
sequences more than quantizing the weights to FP8.
\item \emph{Prediction.} A key's score with a query is $3.0$; the cache
quantization adds an error of $+0.2$. By what factor does that token's
attention weight change before normalization?
\item \emph{Implementation.} Add pre-RoPE key quantization to the lab and
compare it with post-RoPE at 4 bits.
\item \emph{Falsification.} Find a synthetic distribution of keys for which
per-token quantization beats per-channel quantization.
\end{enumerate}
